% ----------------------------------------------------
% Requirements Analysis
% ----------------------------------------------------
\documentclass[class=report,11pt,crop=false]{standalone}
\input{../Style/ChapterStyle.tex}
\begin{document}
% ----------------------------------------------------
\chapter{Feasibility Analysis} \label{ch:FeasibAnalysis}
\vspace{-1cm}
% ----------------------------------------------------
% =====================================================
\section{Traceability Analysis}
Show how the requirements, specifications and testing procedures all link, \autoref{tab:matrix} is provided.

\begin{table}[h]
    \centering
    \caption{Requirements Traceability Matrix}
    \label{tab:matrix}
    \begin{tabular}{|c|c|c|c|}
        \hline
        \# & Requirements & Specifications  & Acceptance Test\\
        \hline
         1 & UR01 &  SP01 & AT01 \\
         \hline
         2 & FR01 & SP02,SP05 & AT03 \\
         \hline
         & &  & \\
         \hline
    \end{tabular}
\end{table}

\subsection{Traceability Analysis 1}
% Explain your thinking around how everything links.

\subsection{Traceability Analysis 2}
% Explain your thinking around how everything links.

% =====================================================
\section{Requirements}
The requirements for a micromouse power module are described in \autoref{tab:requirements}.
\begin{table}[h]
    \centering
    \caption{User and functional requirements of the power subsystem.}    \label{tab:requirements}
    \begin{tabular}{|c|l|}
        \hline
        \textbf{ Requirement ID} & \textbf{Description} \\
        \hline
         UR01 & \\
         \hline
         FR01 & \\
         \hline
         &  \\
         \hline
         &  \\
         \hline
    \end{tabular}
\end{table}

% =====================================================
\section{Specifications}
The specifications, refined from the requirements in \autoref{tab:requirements}, for the micromouse power module are described in \autoref{tab:specifications}.
\begin{table}[h]
    \centering
    \caption{Specifications of the sensing subsystem derived from the requirements in \autoref{tab:requirements}.}    \label{tab:specifications}
    \begin{tabular}{|c|l|}
        \hline
        \textbf{Specification ID} & \textbf{Description} \\
        \hline
         SP01 & User \\
         \hline
         &  \\
         \hline
         &  \\
         \hline
    \end{tabular}
\end{table}

% =====================================================
\section{Acceptance Testing Procedures}
% The table given below is an example of how you can present your ATPs. You are welcome to change or adapt it to your need. The example given is not a valid ATP description.

\begin{table}[h]
  \begin{center}
    \caption{Subsystem acceptance tests}
    \label{tab:tests}
    \begin{tabular}{| >{\centering\arraybackslash}m{1cm} | m{4cm}| m{5cm}| m{5cm}|}
      \hline
      \textbf{ID} & \textbf{Description} & \textbf{Testing Procedure}& \textbf{Pass/Fail Criteria} \\   
      \hline
      AT01 & Powers on & \tabitem \newline\indent\tabitem & \\
      \hline
       &  &\tabitem \newline\indent\tabitem  & \\
      \hline 
    \end{tabular}
  \end{center}
\end{table}



% ----------------------------------------------------
\ifstandalone
\bibliography{../Bibliography/References.bib}
\fi
\end{document}
% ----------------------------------------------------